\chapter{Spécialisation QLoRA et HMoE agentique}
\label{chap:science}

\section{Une mission centrée sur la spécialisation de l'agent}
\label{sec:hmoe-probleme}

La mission scientifique principale du stage a consisté à étudier une
spécialisation frugale de LANCE. Un \gls{llm} généraliste peut expliquer une
vulnérabilité, mais l'agent attendu doit enchaîner des actions beaucoup plus
contraintes : choisir un outil autorisé, produire des arguments conformes à son
schéma, exploiter les observations obtenues et enregistrer le livrable propre à
la phase courante. Ces comportements ne sollicitent pas exactement les mêmes
compétences lors de la reconnaissance, de l'analyse de vulnérabilités, de
l'exploitation ou de la synthèse. Les inscrire uniquement dans le prompt
allonge le contexte et ne garantit pas leur application ; entraîner quatre
modèles complets multiplierait au contraire les besoins en mémoire et en
stockage.

L'hypothèse étudiée est donc la suivante : 

\begin{displayquote}
  un modèle de base quantifié, partagé par plusieurs adaptateurs spécialisés et
  associé à un routage explicite par phase, peut améliorer la conformité
  opérationnelle de l'agent à coût d'entraînement contenu.
\end{displayquote}

Cette hypothèse porte sur le système agentique complet. Elle ne se réduit ni à
la convergence de l'entraînement, ni à la capacité du modèle à compléter une
trace déjà connue. Elle suppose que l'expert sélectionné généralise à des
topologies nouvelles et conserve la capacité d'appeler les outils au bon
moment. Le travail a ainsi couvert toute la chaîne : constitution des corpus,
adaptation \gls{qlora}, service multi-adaptateurs, routage, garde-fous de
contexte et préparation d'un protocole d'évaluation.

\section{Adaptation frugale par QLoRA}
\label{sec:qlora-methode}

La méthode \gls{qlora} conserve les poids du modèle de base gelés et les charge
en quatre bits. Les gradients traversent cette représentation, mais seuls des
adaptateurs de faible rang sont mis à jour. NF4, la double quantification et les
optimiseurs paginés réduisent notamment l'empreinte mémoire de l'apprentissage
\cite{dettmers2023qlora}. Pour une projection de poids \(W\), l'expert \(e\)
apprend la correction

\[
  h = Q_{4}(W)x + \frac{\alpha}{r} B_e A_e x,
\]

où \(Q_{4}(W)\) désigne le poids de base quantifié, \(A_e\) et \(B_e\) les deux
matrices entraînables et \(r\) leur rang. La configuration retenue utilise
\(r=16\), \(\alpha=32\) et un dropout de \(0{,}05\). Les corrections sont
appliquées aux projections de l'attention ainsi qu'aux projections montante,
descendante et de porte du bloc feed-forward. Le calcul s'effectue en
\texttt{bfloat16}, avec accumulation de gradient et activation du
checkpointing de gradient afin de rester compatible avec une carte graphique
de 16\,Gio.

Chaque exemple du corpus de l'expert \(e\) est une conversation
\((x^{(i)},y^{(i)})\in\mathcal{D}_e\), comprenant le contexte, les outils exposés
et la réponse attendue. Un masque \(m_t\) exclut les messages utilisateur et les
résultats d'outils : l'objectif porte uniquement sur les jetons produits par
l'assistant. La fonction minimisée est ainsi

\[
  \mathcal{L}_e =
  -\frac{1}{\sum_{i,t}m_t^{(i)}}
  \sum_{i,t}m_t^{(i)}
  \log p_{\theta,A_e,B_e}
  \left(y_t^{(i)}\mid x^{(i)},y_{<t}^{(i)}\right).
\]

Cette formulation évite d'apprendre à prédire les entrées du protocole et
concentre l'adaptation sur les décisions et productions de l'agent. Une
séparation déterministe conserve 2\,\% de chaque corpus pour l'évaluation. Les
longueurs maximales sont bornées à 6\,144 jetons, sauf pour l'expert de
vulnérabilités limité à 4\,096 jetons. Le prétraitement découpe les traces trop
longues et conserve les outils pertinents, car entraîner sur un contexte plus
large que celui servi créerait un décalage entre apprentissage et inférence.

\section{HMoE agentique : un routage de rôles, pas de jetons}
\label{sec:hmoe-routage}

Le terme \gls{hmoe} désigne ici une architecture agentique hétérogène. Il faut
la distinguer d'un mélange d'experts neuronal classique. Dans ce dernier, une
porte apprise sélectionne, pour chaque jeton et à l'intérieur d'une couche, un
sous-ensemble de blocs feed-forward ; le calcul conditionnel augmente alors la
capacité du réseau sans activer tous ses paramètres \cite{shazeer2017moe}. Dans
LANCE, le modèle neuronal reste unique : le routage intervient au niveau d'une
requête complète et choisit un adaptateur \gls{peft}. Il ne produit donc ni
pondération douce entre experts, ni décision différente pour chaque jeton.

Le routeur exploite d'abord l'identifiant de phase inscrit dans le message
système. Les phases~1 et~6 sont confiées au secrétaire, la phase~2 à la
reconnaissance, la phase~3 à l'analyse de vulnérabilités, et les phases~4 et~5 à
l'exploitation. En l'absence de numéro explicite, des marqueurs de rôle
stabilisent la sélection ; si aucun motif n'est reconnu, le modèle de base sert
de repli. Les identifiants directs des experts restent exposés pour le
diagnostic, tandis que l'identifiant commun \texttt{lance-moe} réalise le
routage automatique. La \cref{tab:hmoe-experts} résume cette décomposition.

\begin{table}[H]
  \centering
  \caption{Décomposition fonctionnelle du \gls{hmoe} agentique}
  \label{tab:hmoe-experts}
  \begin{tabular}{lp{0.16\textwidth}p{0.48\textwidth}}
    \toprule
    Expert & Phases & Compétence principalement apprise \\
    \midrule
    Secrétaire & 1 et 6 & Interprétation de la topologie, continuité des livrables et synthèse finale \\
    Reconnaissance & 2 & Planification des scans, appels d'outils et consolidation des services observés \\
    Vulnérabilités & 3 & Association entre preuves, services et vulnérabilités sans inventer de constat \\
    Exploitation & 4 et 5 & Vérification contrôlée, conservation des preuves et progression vers l'intrusion \\
    \bottomrule
  \end{tabular}
\end{table}

Cette architecture rend la décision de routage interprétable et reproductible.
Elle permet aussi de remplacer un expert sans réentraîner les autres. En
contrepartie, la frontière entre phases est imposée : une requête ambiguë ne
combine pas les compétences de deux adaptateurs. Le changement d'adaptateur
actif est en outre protégé par un verrou de génération ; plusieurs exécutions
concurrentes sont sérialisées afin d'éviter qu'une requête poursuive sa
génération avec les poids d'un autre expert. Des budgets de contexte et de
génération propres aux rôles bornent enfin le coût, mais peuvent tronquer une
trace utile. Ces compromis relèvent de l'architecture agentique et ne sont pas
visibles dans la seule perte d'apprentissage.

\section{Données et observations provisoires}
\label{sec:hmoe-donnees}

Les quatre corpus préparés proviennent de traces structurées selon le protocole
de LANCE. Ils ont été séparés par rôle afin qu'un adaptateur n'apprenne pas des
contrats de sortie incompatibles. Le prévol vérifie le modèle préparé, le rôle
déclaré, les limites de contexte et l'intégrité des fichiers. Les volumes et les
pertes de validation disponibles au moment de la rédaction figurent dans la
\cref{tab:hmoe-apprentissage}.

\begin{table}[H]
  \centering
  \caption{État provisoire des corpus et de l'apprentissage QLoRA}
  \label{tab:hmoe-apprentissage}
  \begin{tabular}{lrrr}
    \toprule
    Expert & Exemples & Contexte maximal & Perte d'évaluation \\
    \midrule
    Secrétaire & 19\,065 & 6\,144 & 0,0421 \\
    Reconnaissance & 9\,257 & 6\,144 & 0,2102 \\
    Vulnérabilités & 6\,951 & 4\,096 & 0,00475 \\
    Exploitation & 14\,629 & 6\,144 & 0,00836 \\
    \bottomrule
  \end{tabular}
\end{table}

Ces valeurs montrent seulement que la vraisemblance des réponses attendues est
mesurable sur la partition réservée. Elles ne sont pas directement comparables
entre experts : les distributions de longueurs, la répétitivité des formats et
la diversité des sorties diffèrent. Une faible perte peut notamment provenir
d'un livrable très structuré et prévisible. Elle ne prouve ni que l'outil choisi
est pertinent, ni que son argument est valide, ni que le constat final est
correct. La perte plus élevée de la reconnaissance peut, inversement, refléter
une plus grande diversité de plans et d'observations sans établir une moindre
efficacité opérationnelle. Ces résultats restent donc des contrôles de
convergence, pas une validation du \gls{hmoe} agentique.

\section{Protocole d'évaluation et ablations}
\label{sec:hmoe-evaluation}

La validation de l'hypothèse de la \cref{sec:hmoe-probleme} doit comparer des
systèmes soumis aux mêmes scénarios, prompts, outils, budgets de génération et
graines. Le jeu d'évaluation doit être distinct des traces d'entraînement et
inclure des topologies ou variantes non observées. Quatre conditions isolent
les apports : le modèle de base avec prompt seul ; un adaptateur QLoRA global
entraîné sur le corpus réuni ; les quatre experts avec sélection forcée de
l'expert correct ; enfin le \gls{hmoe} avec routage automatique. Une ablation
supplémentaire, qui permute volontairement certains experts, mesure si le gain
provient bien de la spécialisation plutôt que du seul ajout de paramètres.

L'évaluation doit être conduite sur deux niveaux complémentaires :

\begin{description}
  \item[Conformité agentique] taux d'appels d'outils syntaxiquement valides,
    respect des outils autorisés, achèvement des phases, conformité des
    livrables et fréquence des récupérations après erreur ;
  \item[Efficacité de l'audit] précision, rappel et score F1 des constats,
    proportion de preuves vérifiées, couverture des chemins d'attaque et coût
    en jetons, temps et mémoire.
\end{description}

Les décisions sémantiques peuvent être assistées par un juge \gls{llm}, mais
celui-ci ne doit pas être l'unique arbitre : les biais de position, de verbosité
et d'auto-préférence sont documentés \cite{zheng2023judge}. Les validations de
schéma, la correspondance avec la \gls{veriteterrain} et le calcul des métriques
doivent rester déterministes lorsque cela est possible. Il faudra enfin publier
les intervalles de confiance sur plusieurs exécutions et journaliser la version
du modèle de base, des adaptateurs, du routeur, des corpus et du code. Ce
protocole permettra de conclure séparément sur trois questions : QLoRA
apporte-t-il un gain par rapport au prompt seul, la séparation des experts
apporte-t-elle un gain par rapport à un adaptateur global, et le routeur
automatique conserve-t-il ce gain sans connaître à l'avance l'expert correct ?
